\normchapter{Fonctions à plusieurs variables}

En physique-chimie, il est rare que l'on puisse décrire un système avec une seule variable. En mécanique, on sait que l'énergie mécanique d'une particule est décrite par ses coordonnées spatiales et ses vitesses, soit 6 variables $(x,y,z,v_x, v_y, v_z)$. Il en est de même en thermodynamique, l'énergie interne par exemple est décrite canoniquement par trois variables $(S, V, N)$ (entropie, volume et nombre de particules).

Pour décrire ces systèmes, il est bien d'avoir quelques notions mathématiques sur les fonctions à plusieurs variables.

\section{Rappel sur la notation différentielle}

\definition{
Soit $f$ une fonction définie et continue au voisinage d'un point $x_0$. Cette fonction admet une dérivée, notée $f'$, telle que : 

\begin{align*}
f'(x_0) = \lim_{\Delta x \to 0} \frac{f(x_0 + \Delta x) - f(x_0) }{\Delta x}
\end{align*}

Si $f$ admet cette limite en $x_0$, on dit que $f$ est dérivable.}

En chimie, on note la dérivée différemment, puisqu'on écrit : 
\begin{equation}
f'(x) = \frac{df}{dx} (x)
\end{equation}

En général, la majorité des chimistes ne feront aucune distinction entre écrire $\frac{df}{dx} (x)$ (l'évaluation de la fonction dérivée au point x, c'est donc un nombre !) et $\frac{df}{dx}$ (la fonction dérivée, c'est un objet mathématique). Ce raccourci est fait par confort, car il permet de facilement passer de la dérivée, à la différentielle ! 

\definition{
\setlength{\columnsep}{5pt}
\begin{wrapfigure}{r}{0cm}% r à droite et l pour à gauche
\includegraphics[width=5cm]{Fig03}
\end{wrapfigure}

On appelle différentielle de $f$, la variation de la fonction $f$, notée $df$, entre un point $x$ et un point $x+dx$ infiniment proche. On note alors : 

\begin{align*}
df (x) = \frac{f(x+dx) - f(x) }{dx} = \frac{df}{dx} (x) dx
\end{align*}

On retrouve dans cette notation les idées de la définition mathématique.}


\section{Dérivées partielles}

Dès à présent, on posera $f$ une fonction à deux variables $(x,y)$ et $g$ une fonction à une variable $z$, on notera leur évaluation respective : 
\begin{align*}
x \to f'(x,y) \quad z \to g(z)
\end{align*}

En particulier, on verra en thermodynamique que les variables d'état dérivent des dérivées de fonctions à plusieurs variables. 

\definition{La dérivée partielle de la fonction $f$ par rapport à x au point (x,y) est la dérivée de la fonction d'une seule variable réelle : 
\begin{align*}
x \to f(x, y) \text{ où  y est constant}
\end{align*}

On note la fonction dérivée partielle : 
\begin{align*}
\left( \frac{\partial f}{\partial x} \right)_y \quad \frac{\partial f}{\partial x}
\end{align*}

La fonction $f$ est dites dérivable selon x au point (x,y), si la fonction $f$ est continue selon x au point (x,y). On peut alors écrire : 

\begin{equation*}
\frac{\partial f}{\partial x} (x,y)  = \lim_{\Delta x \to 0} \frac{f(x + \Delta x, y) - f(x, y)}{\Delta x}
\end{equation*}

On peut faire de même pour la dérivée partielle selon y.
}

Chaque dérivée a un sens bien précis : 

\begin{center}
\begin{tabular}{cC{12cm}}
$\left( \frac{\partial f}{\partial x} \right)_y$ & la dérivée partielle de $f$ par rapport à $x$ quand $y$ est constant \\
$\left( \frac{\partial^2 f}{\partial x^2} \right)_y$ & la dérivée partielle seconde de $f$ par rapport à $x$ quand $y$ est constant \\
$\left( \frac{\partial^2 f}{\partial x \partial y} \right)_y$ & la dérivée partielle de $f$ par rapport à $y$ quand $x$ est constant, suivie de la dérivée partielle par rapport à $x$ lorsque $y$ est constant
\end{tabular}
\end{center}

Les dérivées partielle seconde ne sont que la réalisation successive de deux dérivées partielle : 

\begin{align*}
\left( \frac{\partial^2 f}{\partial x^2} \right)_y &= \left( \frac{\partial}{\partial x} \left( \frac{\partial f}{\partial x} \right)_y \right)_y \\
\left( \frac{\partial^2 f}{\partial x \partial y} \right)_y &= \left( \frac{\partial}{\partial x} \left( \frac{\partial f}{\partial y} \right)_x \right)_y
\end{align*}

Au total, pour une fonction à 2 variables, on peut écrire $2$ dérivées premières et $4$ dérivées secondes : 
\begin{align*}
\left( \frac{\partial f}{\partial x} \right)_y & \left( \frac{\partial f}{\partial y} \right)_x & \left( \frac{\partial^2 f}{\partial x^2} \right)_y & \frac{\partial^2 f}{\partial x \partial y} & \frac{\partial^2 f}{\partial y \partial x} & \left( \frac{\partial^2 f}{\partial y^2} \right)_x
\end{align*}

\exercice{
Soit un gaz parfait dans une enceinte fermée. La loi des gaz parfait s'écrit : 
\begin{equation}
P(V,T) = \frac{nRT}{V}
\end{equation}

où $R$ est la constante des gaz parfait.

\Q Calculer les dérivées partielles premières.
%\begin{align*}
%\frac{\partial P}{\partial V} &= -\frac{RT}{V^2} \\
%\frac{\partial P}{\partial T} &= -\frac{R}{V}
%\end{align*}

\Q Calculer les dérivées partielles secondes.

}



\section{Différentielle totale exacte}
\subsection{Définition}

Dans les fonctions univariées, on sait qu'une fonction g(x) admet systématiquement une différentielle dg, à condition que la fonction $g$ soit dérivable : 
\begin{align*}
dg = \frac{dg}{dx} dx
\end{align*}

Ainsi, à chaque forme différentielle, on peut associer une fonction (c.-à-d. qu'une forme différentielle possède plusieurs primitives dont l'une de ces primitives correspond à g).

Dans les fonctions multivariées, il se trouve que tout ça se complique. On peut très bien définir la forme différentielle d'une fonction à deux variables :

\definition{
On appelle forme différentielle, toute expression mathématique d'une quantité $\delta f$, pouvant s'écrire : 
\begin{align*}
\delta f = X(x,y) dx + Y(x,y) dy
\end{align*}
où $X$ et $Y$ sont deux fonctions quelconques de $x$ et $y$.

Il est à noter que $\delta f$ est une petite quantité, qui peut avoir un sens physique, mais elle n'admet pas forcément de fonction primitive.
}

Le principal souci des formes différentielles, c'est qu'il n'existe pas forcément de fonction (de primitive) à une forme différentielle donnée. Par conséquent, on dit qu'une différentielle est totale exacte lorsqu'elle possède une primitive.

%Soit $f$ une fonction 
%$\Delta x = x - a$ et $\Delta y = y - b$

%\begin{align*}
%\Delta f &= f(x,y) - f(a,b) \\
%&= \frac{\partial f}{\partial x} (x,b) \Delta x + \frac{\partial f}{\partial y} (a,y) \Delta y
%\end{align*}

\exercice{Montrer si les fonctions suivantes admettent une forme différentielle et évaluer-là : 
\Q $f(x,y) = xy$
\Q $f(x,y) = (x-a)^2 + (y-b)^2$

~

Montrer si les formes différentielles suivantes sont exactes : 
\Q $\delta h = 6xy dx + 3x^2 dy$
\Q $\delta w = y^2 dx + x^2 dy$
\Q $\delta z = x^2 y^3 dx +3x^2 y^2 dy$}

\definition{La différentielle totale exacte d'une fonction $f$ à deux variables est définie comme : 
\begin{align*}
df = \frac{\partial f}{\partial x} dx + \frac{\partial f}{\partial y} dy
\end{align*}

On indique par $d$, au lieu de $\delta$, la différentielle lorsqu'elle est exacte.

Si une fonction admet une différentielle totale exacte, alors cette fonction est appelée \textbf{fonction d'état}.}

Une conséquence importante est que lorsque l'on intègre une fonction d'état entre un point initial et un point final, la valeur de son intégrale \textbf{ne dépend pas} du chemin suivi !

\subsection{Conséquence sur les différentielles croisées}

\definition{Une fonction d'état admet une égalité, appelée \textbf{égalité de Schwartz}, entre ses dérivées partielles croisées : 
\begin{align*}
\frac{\partial^2 f}{\partial x \partial y} = \frac{\partial^2 f}{\partial y \partial x}
\end{align*}}

Par conséquent, la réciproque vous indique que si une fonction suit l'égalité de Schwartz, alors cette fonction est une fonction d'état ! 

\exercice{\'{E}crire l'égalité de Schwartz des différentielles totales exactes suivantes : 
\Q $dU (S, V) = T(S,V) dS - P(S, V) dV$
\Q $dH (S, P) = T(S,P) dS + V(S,P) dP$ 
\Q $dG (T, P) = -S(T,P) dT + V(T,P) dP$

~

Montrer si les formes différentielles suivantes sont totales exactes. Conclure si la valeur de l'intégration dépend du chemin suivi.
\Q $\delta f = 6 x y^3 dx + 9x^2y^2 dy$
\Q $\delta g = x^2dx + xydy $}

\subsection{Intégration}
En thermodynamique, on est souvent emmené à décrire des transformations plutôt que des états, c'est-à-dire que mathématiquement, on exprime la variation d'une grandeur au lieu de sa valeur.

Dans les fonctions multivariées, toute la difficulté d'intégrer est inhérente au fait qu'il existe une infinité de chemins d'intégration. 

\begin{center}
\includegraphics[width=0.8\textwidth]{Fig02}
\end{center}

En effet, sur la figure ci-dessus représentant une fonction quelconque, on voit qu'il existe plusieurs façons d'intégrer d'un point $A$ à un point $B$. L'intégrale de cette fonction ne dépendra pas du chemin suivi si elle est une fonction d'état.

Cette idée de fonction d'état est simple à comprendre si on s'appuie sur un analogie à la vie courante. Imaginons que vous cherchiez cet hiver à aller à Bourg-Saint-Maurice en partant de Paris. D'un côté, pour faire la route, vous avez besoin de connaître le denivellé et d'un autre côté, vous avez besoin de connaître la distance. 

Peu importe la route que vous ferez, la différence d'altitude entre Bourg-Saint-Maurice et Paris sera toujours la même, vous aurez toujours le même dénivelé positif. Par contre, la distance, elle, dépendra du chemin que vous suivrez. Ainsi, le dénivelé est une fonction d'état, contrairement à la distance.


\definition{
Si f est une fonction d'état, elle admet donc une différentielle totale exacte. L'intégrale de $f$ d'un point $A$ à un point $B$ est unique et donne : 
\begin{equation*}
\int_A^B df = f_B - f_A
\end{equation*}}

Lorsque l'on connait la forme intégrale d'une fonction d'état, on peut trouver la valeur de son intégrale entre deux ponts en évaluant la fonction d'état en ces deux points. 

Malgré tout, il est bien de savoir intégrer une forme différentielle, en particulier, si l'on ne connait pas sa forme intégrale. Prenons par exemple la forme différentielle :
\begin{equation*}
\delta V = 2 \pi r h dr + \pi r^2 dh
\end{equation*}

Cette fonction respecte l'égalité de Schwartz : 
\begin{align*}
\frac{\partial^2 V}{\partial r \partial h} = \frac{\partial^2 V}{\partial h \partial r} = 2 \pi r
\end{align*}

On peut donc conclure que cette forme différentielle est totale exacte, soit :  
\begin{equation*}
d V = 2 \pi r h dr + \pi r^2 dh
\end{equation*}

On sait exprimer les dérivées partielles, prenons celle par rapport à $r$ : 
\begin{align*}
\frac{\partial V}{\partial r} = 2 \pi r h \stackrel{\text{integration}}{\implies} V({\color{blue} r}, h) &= V({\color{red} r_0}, h) + \int_{\color{red} r_0}^{r \color{blue}}  2\pi r' h dr'\\
&= V({\color{red} r_0}, h) + \pi ({\color{blue} r}^2-{\color{red} r_0}^2) h
\end{align*}

Quand on intègre, la constante d'intégration $V({\color{red} r_0}, h)$ n'est plus qu'une fonction de $h$, que l'on peut estimer à l'aide d'une autre dérivée partielle :

\begin{align*}
\frac{\partial V}{\partial h} (r_0, h)  = \pi r_0^2 \stackrel{\text{integration}}{\implies} V(r_0, h) = V(r_0, h_0) + \pi r_0^2 (h - h_0)
\end{align*}

On peut alors en déduire que : 
\begin{align*}
V (r, h) = \pi (r^2-r_0^2) h + \pi r_0^2 (h - h_0) + V(r_0, h_0)
\end{align*}

En particulier, pour $(r_0, h_0) = 0$, on retrouve $V(r,h) = \pi r^2 h$, ce qui correspond au volume d'un cylindre !  

\exercice{Soit les fonctions d'état suivantes. Intégrez-les de deux façons différentes entre le repère (donc le point (0,0) ou (0,0,0)) et un point M (a,b) ou M (a,b,c) où $(a,b,c)$ sont trois réels positifs.

\Q $S(r, h) = 2 \pi r h + 2 \pi r^2$
 
\Q $V(x,y,z) = x y z$}



\subsection{Recherche des extremums}

\setlength{\columnsep}{1pt}
\begin{wrapfigure}{r}{0cm}% r à droite et l pour à gauche
\includegraphics[width=5.0cm]{Fig01}
\captionof{figure}{Fonction $f(x,y) = (x-x_0)^2 + (y-y_0)^2$ près de son minimum}
\end{wrapfigure}

Dans une fonction à une variable, la ou les racines de sa dérivée nous donne ses extremums (maximum, minimum). Ainsi, si $g$ est dérivable, alors : 
\begin{align*}
\frac{dg}{dx} (x=x^\star) = 0 \implies (x^\star, g(x^\star))\text{ est la coordonnée d'un extremum de g}
\end{align*}

De la même manière, dans une fonction à plusieurs variables, partiellement dérivable selon x et selon y, on a : 
\begin{align*}
df (x^\star, y^\star) = \frac{\partial f}{\partial x} (x^\star, y^\star) \cdot dx + \frac{\partial f}{\partial y} (x^\star, y^\star) \cdot dy = 0 \implies (x^\star, y^\star)\text{ est un extremum de f}
\end{align*}

Dans les derivées à plusieurs variables, un extremum peut être un maximum, un minimum ou un point-selle.

\definition{
Soit $f$ une fonction partiellement dérivable selon x et selon y. Si x, y sont \textbf{indépendants} alors : 

\begin{align*}
df(x^\star, y^\star) = 0 \implies \frac{\partial f}{\partial x} (x^\star, y^\star) = 0 \textbf{ et } \frac{\partial f}{\partial y} (x^\star, y^\star) = 0
\end{align*}

Pour rechercher les extremums d'une telle fonction, il faut que les deux dérivées partielles soient nulles au point étudié !
}

Vous ne pouvez pas conclure sur le fait qu'un point soit un extremum sans étudier toutes les dérivées premières d'une fonction multivariée.